% CoDIT regular paper (up to 6 pages, IEEE conference double-column). Build: pdflatex, bibtex, pdflatex x2.
% Every number below is traced in paper/codit2026/README.md; \tbd marks a value still to be measured.
% SCOPE (researcher, 2026-11-07): the PFB as a safety filter for CLASSICAL local planners only; no learned
% planner is mentioned. PRELIMINARY: all results are tune-set results; get the supervisor's approval first.
\documentclass[conference]{IEEEtran}
\IEEEoverridecommandlockouts
\usepackage{amsmath,amssymb,amsfonts}
\usepackage{graphicx}
\usepackage{booktabs}
\usepackage{xcolor}
\usepackage{cite}
\usepackage{url}
\newcommand{\tbd}[1]{\textcolor{red}{[#1]}}

\begin{document}

\title{Protective-Field Barriers: A LiDAR Control Barrier Function Matched to Industrial Safety Fields for Classical Local Planners}

\author{\IEEEauthorblockN{\tbd{Author 1}, \tbd{Author 2}}
\IEEEauthorblockA{\tbd{Affiliation, City, Country}\\
\tbd{email}}}

\maketitle

\begin{abstract}
Mobile robots in industrial facilities run above a certified safety layer: safety laser scanners whose
speed-dependent protective fields trigger a latched protective stop whenever an object enters them. A
local planner that does not model these fields triggers stops routinely, at corners, in narrow aisles and
near unmapped obstacles. We propose the Protective-Field Barrier (PFB), a discrete-time control barrier
function computed directly on the LiDAR scan, whose safe set is the complement of the certified layer's
own trigger set: a continuous, speed- and direction-dependent field built from the same configuration as
the safety function. The PFB filters the command of any local planner through a small quadratic
program. On a holonomic platform in simulated warehouse aisles, adding the PFB to a tuned Regulated Pure
Pursuit planner reduces protective stops from 86.1 to 0.59 per 100\,m driven and the time spent latched
from 42\,\% to 0.05\,\%, without lowering its driving speed. Its remaining failures come from the
planner's own speed regulation and from routes that no route-following controller can pass. A tuned
field-aware model predictive path integral controller reaches a comparable stop rate. These results are
preliminary, from the tuning set of a preregistered study.
\end{abstract}

\begin{IEEEkeywords}
control barrier functions, safety filter, local planning, protective stop, mobile robots
\end{IEEEkeywords}

\section{Introduction}
Autonomous mobile robots in warehouses share space with people and equipment, and the standards that
govern them require a certified safety function independent of the navigation software
\cite{iso3691-4,iso13849-1}. In practice this is a safety laser scanner \cite{iec61496-3} that monitors
protective fields sized to the robot's speed and stops the robot, with a latch, when anything enters the
active field. The navigation stack above it---a global planner and a local planner that tracks the
route---is not certified and may choose any command. The two layers do not share a model: a local
planner that treats obstacles through a costmap or a fixed inflation radius enters the fields at
corners, in narrow aisles and near unmapped obstacles, and every entry costs a stop, a restart and often a
manual intervention. Protective stops per distance travelled are therefore a practical measure of how
well a planner and its safety layer fit together.

Control barrier functions (CBFs) \cite{ames2017cbfqp,ames2019cbf} keep a system inside a safe set: a
function $h(x)$, nonnegative on the safe set, defines a constraint on the input that keeps it so, and a
quadratic program (QP) finds the input closest to a nominal command that satisfies it. As a safety
filter, a CBF can sit beneath an existing planner without changing it \cite{wang2022ensuring}. Most
LiDAR- or map-based CBFs for mobile robots, however, encode a geometric clearance---a distance margin to
obstacles \cite{raja2024ogmcbf,zeng2025relvel}---with no defined relation to the certified field; such a
filter can be too conservative, or still trigger stops.

This paper makes the barrier's safe set the complement of the certified layer's trigger set. Our
contributions are:
\begin{itemize}
\item the PFB, a discrete-time CBF on the sector scan whose zero set is a continuous, speed- and
direction-dependent protective field built from the certified configuration, and its filter as a
per-regime QP that respects the platform's actuation and acceleration limits;
\item an evaluation in which every planner runs above the same emulated certified layer, with
protective stops per 100\,m as the primary metric, and the share of time latched and the moving speed
as diagnostics;
\item preliminary results on a holonomic platform: the PFB removes nearly all protective stops from a
classical planner without slowing it, with an analysis of the failures that remain.
\end{itemize}

\section{Related Work}
\subsection{Classical local planners}
Pure Pursuit \cite{coulter1992purepursuit} steers toward a lookahead point on the route. Regulated Pure
Pursuit (RPP) \cite{macenski2023rpp}, the default path tracker of Nav2 \cite{macenski2020nav2}, scales
the speed down with path curvature and with proximity to obstacles. The Dynamic Window Approach (DWA)
\cite{fox1997dwa} samples admissible velocities within one acceleration window and scores short rollouts
against the route, the heading and the clearance. Model Predictive Path Integral control (MPPI)
\cite{williams2017mppi} optimises a sampled control sequence by importance weighting of rollout costs and
accepts arbitrary costs, including one on entering a safety field. All of them reason about obstacles
through the planner's own clearance or cost; none is given the certified layer's fields as a hard
constraint.

\subsection{Safety layers and barrier filters}
Shielding \cite{alshiekh2018shielding} and CBF-based filters
\cite{cheng2019endtoend,wang2022ensuring,knoedler2025onthefly,guo2025radar} correct each command of an
upstream controller to keep the state in a safe set. For mobile robots, the safe set is usually a
clearance from obstacles in a map or a scan \cite{raja2024ogmcbf}, sometimes enlarged with the relative
velocity \cite{zeng2025relvel}. The PFB belongs to this family; it differs in what it protects: not a
clearance chosen by the designer, but exactly the region whose entry the certified safety function
punishes with a stop, at the speed and direction the robot will have.

\section{Problem Formulation and System Design}
\subsection{Platform and certified layer}
The robot is a holonomic (mecanum) base commanded by a body-frame twist $u=(v_x,v_y,\omega)$ at
$\Delta t=0.05$\,s, within an axis box ($|v_x|,|v_y|\le0.5$\,m/s, $|\omega|\le1.5$\,rad/s), a wheel
rim-speed budget and an acceleration limiter. Two S300 safety laser scanners give a 120-bin sector scan
with a validity flag per bin; an unobserved bin is never treated as free. The certified layer is
emulated independently of the filter (the two share configuration, never code): it selects a protective
field from discrete speed bands, triggers a latched stop when a valid return lies inside the field in
force, and releases the latch only under a restart condition. Collisions are decided by the simulator's
contact check alone.

\subsection{The protective field}
For travel at speed $V$ in direction $\theta$ (body frame) with yaw rate $\omega$, the field's boundary
distance along bearing $\psi$ is $\rho(\psi;V,\theta,\omega)$, obtained from a support function over
normals $n$. With $c=\cos(n-\theta)$ and $\ell=\max(0,c)$, the support is
\begin{equation}
\begin{aligned}
\sigma(n)=\max\Big(&S\big(n;\pm(E(|\omega|)+E(0))\big)\\
&+\max\big(D(V)\,\ell,\,s(c,V)\big)\\
&+K(V,\omega)+\max\big((m-m_\perp)\ell,R\big)+m_\perp,\\
&S\big(n;\pm E(|\omega|)\big)+R+m_\perp\Big),
\end{aligned}
\end{equation}
where $S(n;\pm e)$ is the footprint's support swept over headings in $[-e,e]$, $E(\omega)$ the heading
change reachable within the reach time $T$ (the certified response time plus one control step,
$T=0.15$\,s) and braking, $D(V)=VT+V^2/(2a_b)$ the travel over the reach time and braking, $s$ the
displacement the limiter allows off the travel direction, $K$ the drift of a body turning while it
translates, $R$ the at-rest reach, and $m$, $m_\perp$ the margins ahead and abeam. The second term is the
at-rest field, the infimum of the moving fields. The PFB uses the field of the band that the reachable
speed selects, so its trigger set is the certified layer's at the same state.

\subsection{The barrier and the filter}
For each valid return $i$ within reach, the barrier is the return's distance outside the field,
$h_i=r_i-\rho(\psi_i;\cdot)-\delta$, with $\delta$ the measured binning residual. For a candidate
command $u$, the return is propagated in the body frame over the prediction horizon
$\tau=\Delta t+\tau_\ell$ (one step plus a declared lag allowance), and the discrete-time barrier
condition is imposed:
\begin{equation}
\begin{aligned}
g_i(u)&=r_i(u)-\rho\big(\psi_i(u);\theta(u),W(u)\big)-\delta\\
&\ge(1-\alpha\Delta t)\max\big(h_i,0\big),
\end{aligned}
\label{eq:cbf}
\end{equation}
where $\theta(u)$ is the field's heading and $W(u)=\max(|\hat\omega|,|\omega|)$ composes the estimated
and commanded yaw rates. Clipping at zero requires a field that is breached now to be clear after one
step. The filtered command solves, for each admissible regime (the at-rest field while creeping, else a
speed band),
\begin{equation}
\min_u\ (u-u_{\mathrm{nom}})^\top \mathrm{diag}(1,1,\rho_{\mathrm{out}}^2)\,(u-u_{\mathrm{nom}})
\label{eq:qp}
\end{equation}
subject to the rows \eqref{eq:cbf} linearised about a reachable point near $u_{\mathrm{nom}}$, the wheel
polytope, the axis box, the limiter's one-step envelope and the regime's speed range; $\rho_{\mathrm{out}}$
is the footprint's outer radius, which weights yaw against translation. The true rows are checked at
the solution without tolerance, and the QP is re-linearised once if needed. The best feasible regime
wins; if none is feasible, the filter commands maximum braking and counts the step as a fallback. Under
the certified latch, only a command the restart condition could release is passed. The QP is solved with
Clarabel \cite{goulart2024clarabel}.

The planner whose command is filtered is unchanged: the PFB needs only its nominal twist, the scan and
the velocity estimate. Here it filters RPP; the combination is called B7.

\section{Experiments}
\subsection{Setup}
Experiments run in CoppeliaSim \cite{rohmer2013vrep} with the youBot model in generated warehouse halls
of three aisle-width strata (narrow, medium, wide). Each scenario adds one or two unmapped obstacles near
the route, which the global planner never sees and around which it never replans. Every planner receives
the same A* route, runs above the same certified layer, and has the same recovery behaviour after a stop.
Compared: \textbf{B1} Pure Pursuit with its own lookahead distance; \textbf{B2} RPP; \textbf{B3} DWA;
\textbf{B6} field-aware MPPI, whose rollout cost penalises scan points inside the certified field;
\textbf{B7} B2's tuned parameters plus the PFB. B1, B2, B3 and B6 were each tuned with the same budget of
24 candidates on the tuning set (selection: fewest stops among candidates within 5 percentage points of
the best success rate). The PFB's single gain was selected the same way with B7, from
$\alpha\in\{1,2,5,10,20\}$\,s$^{-1}$, giving $\alpha=1$. A privileged oracle tracker, which knows every
obstacle but follows the route, serves as a positive control.

The primary metric is protective stops per 100\,m of planned route; stops per 100\,m driven, the share
of steps latched and the moving speed (mean speed while not latched) are reported beside it. Episodes
last at most 200\,s. All results are on the tuning set (60 scenarios, 633\,m of planned route) and are
preliminary; a preregistered evaluation on held-out layouts is pending.

\subsection{Results}
Table~\ref{tab:main} summarises the planners. Adding the PFB to B2 reduces protective stops from 86.1 to
0.59 per 100\,m driven and the latched share from 42\,\% to 0.05\,\%, and doubles the success rate
(0.25 to 0.50). Table~\ref{tab:strata} shows the effect in every aisle width.

\begin{table}[t]
\centering
\caption{Tuning set, 60 scenarios (preliminary).}
\label{tab:main}
\setlength{\tabcolsep}{3pt}
\begin{tabular}{lccccc}
\toprule
Planner & Success & \multicolumn{2}{c}{Stops /100\,m} & Latched & Moving \\
 & & planned & driven & share & m/s \\
\midrule
B1 Pure Pursuit & 0.017 & \tbd{} & 173.1 & 0.875 & 0.158 \\
B2 RPP & 0.250 & \tbd{} & 86.1 & 0.420 & 0.059 \\
B3 DWA & 0.667 & 20.5 & 23.8 & 0.064 & 0.070 \\
B6 field-aware MPPI & 0.833 & 0.47 & 0.44 & 0.001 & 0.091 \\
\textbf{B7 RPP + PFB} & 0.500 & 0.47 & 0.59 & 0.001 & 0.048 \\
\midrule
Oracle (privileged) & 0.717 & 49.0 & 59.1 & 0.377 & 0.110 \\
\bottomrule
\end{tabular}
\end{table}

\begin{table}[t]
\centering
\caption{B2 and B7 by aisle-width stratum, 20 scenarios each.}
\label{tab:strata}
\setlength{\tabcolsep}{4pt}
\begin{tabular}{lcccc}
\toprule
 & \multicolumn{2}{c}{Success} & \multicolumn{2}{c}{Stops /100\,m driven} \\
Stratum & B2 & B7 & B2 & B7 \\
\midrule
Narrow & 0.30 & 0.40 & 59.4 & 0.0 \\
Medium & 0.15 & 0.40 & 88.7 & 0.0 \\
Wide & 0.30 & 0.70 & 118.7 & 1.76 \\
\bottomrule
\end{tabular}
\end{table}

\emph{The filter does not slow the planner.} On the 13 routes where B7 timed out with no obstacle in
play, B7 and B2 drive at the same speed when not latched (0.057\,m/s pooled for both). B2 spends 27\,\%
of those episodes latched, with 69 stops; B7 none. B7's slowness is B2's own: the tuned RPP begins to
slow at 1.18\,m of clearance, which in 0.7--0.9\,m aisles holds it near its minimum speed throughout;
alone in an empty synthetic aisle it covers 0.73\,m in 10\,s, where Pure Pursuit and DWA cover 4.3 and
4.2\,m.

\emph{Where B7 fails.} B7's 30 failures are all timeouts, 28 of them without a single stop. Seventeen are routes blocked
by an unmapped obstacle across the route, 0.01--0.23\,m from it: the privileged oracle fails exactly
these 17, so no route-following controller passes them. The other 13 are horizon timeouts on long routes
(11.4--17.0\,m), on which B7 covered 68--97\,\% of the route before the 200\,s limit. On the 43 routes a
route follower can traverse, B7 succeeds on 30 (70\,\%).

\emph{The gain.} Table~\ref{tab:alpha} shows the sensitivity to $\alpha$. Success is flat (the 200\,s
horizon caps it), while stops rise with $\alpha$: a smaller $\alpha$ lets the barrier shrink more slowly
per step and keeps the robot further from the field. At $\alpha=1$\,s$^{-1}$ the QP is infeasible on
332 steps, on which the filter brakes at the limiter's bound; a robot at 0.3\,m/s with a wall 0.95\,m
ahead is such a case, filtered smoothly at $\alpha=2$.

\begin{table}[t]
\centering
\caption{B7 against the PFB gain $\alpha$ (60 scenarios).}
\label{tab:alpha}
\setlength{\tabcolsep}{4pt}
\begin{tabular}{lccccc}
\toprule
$\alpha$ (s$^{-1}$) & 1 & 2 & 5 & 10 & 20 \\
\midrule
Success & 0.500 & 0.500 & 0.517 & 0.500 & 0.517 \\
Stops /100\,m driven & 0.59 & 2.94 & 9.63 & 11.98 & 12.93 \\
Fallback steps & 332 & 125 & 74 & 2807 & 8142 \\
\bottomrule
\end{tabular}
\end{table}

\emph{Against the other planners.} Pure Pursuit is latched 83--90\,\% of every episode at every one of
its 24 candidates: it stops at shelf corners and its restart never clears. DWA succeeds more often than
B7 (0.667) but with 23.8 stops per 100\,m. The tuned field-aware MPPI, which is given the certified field
in its rollout cost, is the strongest planner on the tuning set: 50/60 successes and 3 stops, the same
count as B7. With so few events, the tuning set cannot separate B6 and B7 on stops.

\emph{Compute.} Per-step decision time on the host (p95) is 0.02\,ms (B1), 0.03\,ms (B2), 4.4\,ms (B3),
37.1\,ms (B6, 1024 samples $\times$ 20 steps), and \tbd{ms} for the PFB, all within the 50\,ms control
period.

\section{Discussion and Conclusion}
Matching the barrier's safe set to the certified field is enough to remove nearly all protective stops
from an existing classical planner, in every aisle width, at no cost in driving speed and without
changing the planner. The failures that remain belong to the planner and the task, not to the filter:
a conservative speed regulation that the 200\,s limit exposes, and obstacles that no route follower can
pass without replanning. A field-aware MPPI reaches a similar stop rate by building the field into its
cost; the PFB obtains it as a drop-in filter beneath a planner that knows nothing of the field. Which of
the two is preferable on stops, success and computational cost is a question for the confirmatory
evaluation, which the tuning set cannot answer.

The evidence is preliminary: one tuning set of 60 scenarios, few events, one platform and simulator, and
an emulated rather than real safety function. The confirmatory evaluation is preregistered, on held-out
layouts with a test size set by a power analysis, and reports the rate ratio of protective stops per
planned length against each planner. Future work will report it, examine planners re-tuned with the
filter in place, and transfer the filter to hardware with certified scanners.

\bibliographystyle{IEEEtran}
\bibliography{references}

\end{document}
